\documentclass[12pt,a4paper]{report}
\usepackage{setspace} % For line spacing
\usepackage{geometry} % For setting margins
\usepackage{graphicx} % For figures
\usepackage{titlesec} % For section formatting
\usepackage{hyperref} % For references and links
\usepackage{fancyhdr} % For headers and footers
\usepackage{tocloft} % For table of contents customization

% Margins
\geometry{
    left=1.25in,
    right=1in,
    top=0.75in,
    bottom=0.75in
}

% Line spacing
\onehalfspacing % For 1.5 line spacing, or use \doublespacing for double spacing

% Chapter and Section formatting
\titleformat{\chapter}[hang]{\LARGE\bfseries}{\thechapter.}{20pt}{\LARGE\bfseries\centering}
\titleformat{\section}[hang]{\Large\bfseries}{\thesection}{15pt}{\Large\bfseries}
\titleformat{\subsection}[hang]{\large\bfseries}{\thesubsection}{10pt}{\large\bfseries}



% Header/Footer settings
\fancypagestyle{plain}{
  \fancyhf{}
  \fancyhead[L]{2024-2025}
  \fancyhead[R]{AIET,CSE}
  \fancyfoot[C]{\thepage}
  \renewcommand{\headrulewidth}{0.4pt}
}

% Begin document
\begin{document}

% Cover Page
\begin{titlepage}
    

	
	
\centering
\vspace*{-0.75in}

\large 
\textbf{VISVESVARAYA TECHNOLOGICAL UNIVERSITY\\
	Jnana Sangama, Belagavi - 590 018}\\
\vspace{0.2cm}
\begin{figure}[h]
	\centering
	\includegraphics[height=3cm]{vtu.png}
\end{figure}
\small{\textbf{PROJECT REPORT ON}}\\

\LARGE{\textbf{CV- Analyzer}}
\vspace{0.5cm}

\large \textit{Report submitted in partial fulfillment for the Completion of }\\{\textbf{Mini Project}}\\
\ in \\\textbf{Computer Science and Engineering}
\vspace{0.5cm}\\
{\textbf{Submitted by\\}}


\begin{tabular}{lll}
	
Advith Suvarna & \hspace{5cm} & 4AL22CS005\\
Akhifa Sheikh & \hspace{5cm} & 4AL22CS006\\
Akhilesh & \hspace{5cm} & 4AL22CS007\\
Akshay RB & \hspace{5cm} & 4AL22CS008\\
	
\end{tabular}
\vspace{0.5cm}\\
\textit{Under the Guidance of}


\Large{\textbf{Ms. Vidya \hspace{3cm}    \hspace{3cm}Ms.Deeksha.M}}\\
\textit{Senior Associate Professor \hspace{2cm} Associate Professor}\\

\begin{figure}[h]
	\centering
	\includegraphics[height=2.7cm]{aiet-logo.png}
\end{figure}


\begin{center}
	\scriptsize\textbf{DEPARTMENT OF COMPUTER SCIENCE AND ENGINEERING}\\
	\small\textbf{(Accredited by NBA for the Academic years 2019-25)}	
\end{center}
\begin{center}
	\vspace{0.1cm}
	\large\textbf{ALVAS INSTITUTE OF ENGINEERING AND TECHNOLOGY}\\
	\small\textbf{Shobhavana Campus Mijar, Moodbidri,Karnataka 574225\\
2023 - 24}\\
\end{center}
\end{titlepage}

% Abstract
\chapter*{Abstract}
\textbf{\textit{}CV Analyser} is a mini-project that functions as a resume analyzer with Applicant Tracking System (ATS) capabilities, developed in Python and leveraging the Gemini API. This tool allows users to upload resumes, which are then parsed and analyzed based on predefined prompts. By employing the Gemini API, the project automates the extraction of essential details, performs keyword matching, and generates a compatibility score for potential job roles. The output offers users insights into strengths, areas for improvement, and alignment with job requirements, enhancing job seekers' chances of standing out in ATS-filtered application processes.
\newpage
\tableofcontents
\newpage



% Introduction Chapter
\chapter{Introduction}
\textbf{CV Analyzer} In the competitive landscape of job applications, Applicant Tracking Systems (ATS) have become a standard tool for recruiters to efficiently screen, filter, and evaluate a vast number of resumes. While ATS can streamline hiring by identifying the most qualified candidates, they also present a challenge for job seekers who may struggle to meet specific formatting or keyword requirements to pass through initial screenings. CV Analyser is designed to address these challenges by acting as an ATS-friendly resume evaluation tool, allowing users to upload their resumes and receive an analysis based on predefined criteria. \\

The project, developed using Python and integrated with the Gemini API, performs automated parsing and evaluation of resume content to generate valuable insights. Through predefined prompts, the CV Analyser identifies key resume sections—such as experience, education, skills, and achievements—and assesses them based on relevance, keyword density, and alignment with job-specific requirements. By analyzing these attributes, the system calculates a compatibility score, providing users with feedback on areas of strength and suggesting improvements for better ATS optimization.\\
\\
In addition to keyword matching, CV Analyser also checks for structural and formatting factors, which can significantly impact a resume’s ability to pass ATS filters. Through the Gemini API’s robust data-processing capabilities, the project can offer a comprehensive analysis that supports job seekers in optimizing their resumes, making them more appealing not only to ATS but also to hiring professionals. Overall, CV Analyser aims to empower users by offering actionable insights, helping them enhance their resumes to improve their chances of standing out in today's automated recruitment processes.


% Literature Review
\chapter{Literature Review}

\textbf{1. Digital Transformation in Recruitment and ATS Systems}

\begin{itemize} \item \textbf{”Automating Inequality: How High-Tech Tools Profile, Police, and Punish the Poor”} by Virginia Eubanks. This book explores how automated systems are increasingly used in decision-making, including hiring, impacting diverse groups of job seekers. The insights highlight how applicant tracking systems (ATS) can affect the visibility of candidates, particularly in cases where keyword matching and algorithmic evaluations are used. For the CV Analyser, Eubanks' work underscores the importance of developing fair and transparent algorithms to ensure all applicants are given an equitable opportunity.”Automating Inequality: How High-Tech Tools Profile, Police, and Punish the Poor” by Virginia Eubanks. This book explores how automated systems are increasingly used in decision-making, including hiring, impacting diverse groups of job seekers. The insights highlight how applicant tracking systems (ATS) can affect the visibility of candidates, particularly in cases where keyword matching and algorithmic evaluations are used. For the CV Analyser, Eubanks' work underscores the importance of developing fair and transparent algorithms to ensure all applicants are given an equitable opportunity.[1].

\item \textbf{"Talent Wins: The New Playbook for Putting People First”} by Charan, Barton, and Carey. This book provides an in-depth look into talent management strategies and the use of technology in identifying potential hires. The CV Analyser benefits from these insights by emphasizing the importance of data-driven methods to support recruiters and job seekers in improving the quality of resume screening and talent matching. [2]. \end{itemize}

\textbf{2.Resume Optimization and Keyword Matching in ATS}

\begin{itemize} \item \textbf{”The ATS Resume Test: Guide to Landing More Interviews”} by Lewis Lin. Lin’s book discusses ATS resume optimization techniques, with a focus on keyword density, formatting, and structure that can affect the pass rate through ATS filters. These principles are essential for the CV Analyser’s functionality as it aims to provide actionable feedback on resume structure, keywords, and compatibility with ATS. [3].


\item \textbf{”Resumes that Pop! Designs that Reflect Your Personal Brand”} by Pat Criscito. This work emphasizes the importance of creating visually appealing and strategically designed resumes that also align with ATS parsing requirements. Criscito's insights highlight the need for tools like the CV Analyser to support users in balancing resume aesthetics and ATS readability, ensuring both visibility and professionalism in the digital hiring process. [5]. \end{itemize}

\textbf{3.User Experience (UX) in Career Development Tools}

\begin{itemize} \item \textbf{”Designing for the Digital Age”} by Kim Goodwin. Goodwin’s book provides a comprehensive guide to UX design in digital tools, focusing on usability, clarity, and user-centered interactions. These UX principles are crucial for the CV Analyser to ensure intuitive navigation, clear feedback, and a seamless user experience for job seekers looking to improve their resumes.[6].

\item \textbf{”A Project Guide to UX Design: For User Experience Designers in the Field or in the Making”} by Russ Unger and Carolyn Chandler. This work offers strategies for developing accessible, user-friendly digital products that align with user goals. By following Unger and Chandler’s UX design principles, the CV Analyser can ensure an interface that is engaging and supportive, enabling users to receive meaningful, personalized feedback on their resumes. [7]. \end{itemize}

\textbf{4. Trust, Transparency, and Feedback in Automated Career Tools}

\begin{itemize} \item \textbf{”The Trust Factor: The Science of Creating High-Performance Companies”} by Paul J. Zak. Zak’s work on trust-building in digital systems examines how transparent feedback and actionable insights foster user confidence. In the context of CV Analyser, implementing clear and understandable feedback mechanisms will help users feel confident in the reliability of the ATS analysis provided, encouraging constructive use. [8].

\item \textbf{”The Feedback Loop: Using Feedback to Drive Performance”} by Benedict Bensimon. Bensimon’s book covers the importance of constructive feedback in personal and professional growth. The CV Analyser uses Bensimon’s concepts by delivering specific, actionable advice on resume improvements, thereby helping users understand their strengths and areas of improvement in alignment with ATS requirements. [9].

\item \textbf{"PCI Compliance: Understand and Implement Effective PCI Data Security Standard Compliance"} by Branden Williams and Anton Chuvakin. For any platform dealing with online payments, PCI compliance is essential. This book offers guidance on meeting security standards, which is crucial for CynTM’s secure payment gateway and handling of sensitive user data [10]. \end{itemize}

\textbf{5. Ethics and Bias in Automated Screening Tools}

\begin{itemize} \item \textbf{”Weapons of Math Destruction: How Big Data Increases Inequality and Threatens Democracy”} by Cathy O'Neil. This book discusses the ethical implications of algorithm-driven assessments, especially in hiring practices. O'Neil's perspectives are relevant to the CV Analyser as it strives to develop ethical algorithms that avoid bias in resume scoring and ensure equal chances for diverse applicants.[11].

\item \textbf{”AI Ethics”} by Mark Coeckelbergh. Coeckelbergh examines ethical frameworks and best practices for AI applications, emphasizing fairness, accountability, and transparency. This literature is applicable to CV Analyser’s development, as it highlights the importance of incorporating ethical standards to provide fair resume assessments without amplifying biases inherent in some ATS systems. [12]. \end{itemize}
\chapter{Existing System}

\begin{enumerate} \item \textbf{Jobscan:}\\Function: Jobscan is an online tool that allows users to upload their resumes and compare them against a specific job description. Using keyword-matching technology, it evaluates the resume’s relevance to the job and provides a compatibility score.\\
How It Works: Jobscan analyzes keywords, phrasing, and formatting, offering users a percentage match score based on its proprietary ATS algorithm. It also identifies missing keywords or sections and provides suggestions for improving resume structure.\\
Limitations: Although Jobscan provides helpful feedback, its suggestions can be generalized, and it lacks customization for different ATS requirements. Additionally, it focuses mainly on keyword alignment rather than on structural elements that improve overall readability and compatibility across diverse ATS systems.

\item \textbf{Resume Worded:}\\Function: Resume Worded provides automated feedback on resume content, structure, and keyword alignment based on job description input. It also offers an "ATS Optimization" score, focusing on ATS-friendliness and key phrases relevant to the job.\\
How It Works: The platform uses AI to scan for critical ATS factors, such as keyword density, section titles, and readability, giving users insights into areas like length, formatting, and ATS compliance. It also suggests improvements to meet ATS standards.\\
Limitations: Resume Worded can sometimes produce overly generic feedback that may not apply across different industries. It also focuses primarily on ATS keyword matching rather than personalized advice for specific job roles, which can limit its usefulness for specialized or technical fields.

\item \textbf{Enhancv:}\\ Function: Enhancv combines resume building and ATS scoring in one platform, guiding users through template selection and providing real-time suggestions to ensure ATS compatibility.\\
How It Works: Enhancv uses predefined templates designed to meet ATS standards while also allowing for custom adjustments. As users build their resumes, Enhancv provides a “fit score” and offers advice on wording, structure, and layout for ATS optimization.\\
Limitations: While Enhancv excels in ease of use and aesthetic options, its limited customization for specific ATS algorithms may result in less effective parsing by some ATS platforms. Additionally, its feedback can be superficial, with less focus on advanced optimization metrics like language relevancy or job-specific tailoring.

\item \textbf{SkillSyncer:} Function: SkillSyncer compares a user’s resume to a specific job posting, highlighting matching skills, keywords, and qualifications, and providing an ATS compatibility score.\\
How It Works: SkillSyncer scans resumes for industry-specific keywords and qualifications, aligning them with a target job posting. The tool identifies missing key phrases, skill gaps, and section order recommendations to enhance ATS compliance.\\
Limitations: SkillSyncer is heavily focused on keyword matching, which, while useful, doesn’t address more complex elements of ATS optimization like structure, readability, or dynamic scoring for continuous improvement.
\end{enumerate}
% Proposed System
\chapter{Proposed System}


\section{Features of the Proposed System}

\begin{itemize} \item \textbf{User-Friendly Interface:} CV Analyser features a clean, easy-to-navigate interface where users can upload resumes, input job descriptions, and view analysis results quickly. The layout is intuitive, making it easy for users to manage resume uploads and view specific areas for improvement, ensuring a smooth user experience.



\item \textbf{Comprehensive Resume Analysis:}The platform offers a detailed analysis covering essential ATS requirements, including keyword relevance, formatting, structure, and overall readability. This analysis allows users to align their resumes with specific job descriptions, ensuring maximum compatibility and effectiveness.
\item \textbf{ATS Compatibility Scoring:}  CV Analyser includes a scoring system that evaluates ATS-friendliness, ranking resumes based on factors such as keyword presence, section structure, and formatting consistency. The score helps users understand how their resumes measure up to ATS standards and provides actionable feedback for improvement.
\item \textbf{
Targeted Keyword and Phrase Suggestions:} To increase resume relevance, CV Analyser suggests industry-specific keywords and phrases based on the job description provided by the user. This feature enhances the resume’s effectiveness by aligning it more closely with hiring managers' expectations and ATS algorithms.

\item \textbf{Real-Time Feedback and Update System:}As users make adjustments to their resumes, they receive real-time feedback and an updated compatibility score, helping them track the impact of each change immediately. This iterative process allows for faster, more efficient resume optimization.
\item \textbf{Secure Data Handling and Privacy Protection:} The platform ensures that user data, including resumes and job descriptions, are stored securely with advanced encryption standards. This safeguards users' personal information and provides a safe environment for managing sensitive career-related documents.

\item \textbf{Dashboard for Resume Management:} Users have access to a personalized dashboard where they can upload, organize, and track different resume versions. This feature is particularly useful for users who need to maintain multiple resumes tailored for different job applications, making it easy to manage, update, and monitor their progress.

\item \textbf{Responsive Design for Multi-Device Access:}  Built with responsive design principles, CV Analyser is accessible on desktops, tablets, and mobile devices, enabling users to optimize their resumes on the go and check their scores from any device.
\end{itemize}
% Problem Statement
\chapter{Problem Statement} In today’s competitive job market, job seekers struggle to create resumes that are compatible with Applicant Tracking Systems (ATS), which many companies use to screen candidates. Traditional resume formats often fail to meet ATS requirements, resulting in qualified candidates being overlooked. Current resume analysis tools provide limited, generalized feedback that doesn’t cater to industry-specific roles or adapt to the nuances of different ATS algorithms. Therefore, there is a need for a tool that empowers job seekers by offering targeted, real-time feedback to enhance their resumes for ATS compatibility and maximize their chances of landing interviews.

\section{Objectives} The proposed solution aims to accomplish the following objectives:

\begin{itemize} \item \textbf{To develop a user-friendly CV Analyser platform} that allows users to upload resumes, receive detailed ATS-focused analysis, and make data-driven improvements based on real-time scoring. \item \textbf{To provide} industry-specific keyword and phrase recommendations tailored to job descriptions, ensuring resumes meet specific ATS criteria.\item \textbf{To implement} a comprehensive scoring system that evaluates resumes for factors such as formatting, keyword alignment, and readability, helping users understand how ATS-friendly their resumes are. \item \textbf{To enable} secure storage and management of multiple resume versions within a personalized dashboard, allowing users to track and refine resumes across job applications. \end{itemize}
\newpage

\section{Methodology}
\begin{itemize} \item \textbf{Platform Development and User Interface Design:} Develop an intuitive user interface that facilitates easy resume uploads and displays analysis results clearly, focusing on user accessibility and smooth navigation across devices.


\item \textbf{Resume Parsing and ATS Scoring:}Implement resume parsing algorithms to analyze key sections, format structure, and keyword density. Create an ATS scoring model that evaluates resumes based on industry-standard ATS requirements, including keyword matching, structure alignment, and format consistency.
\item \textbf{Keyword and Phrase Suggestion Module:} Develop a module that extracts job-specific keywords and phrases from user-inputted job descriptions, offering tailored recommendations to enhance resume relevance for target roles.

\item \textbf{Real-Time Feedback Integration:} Incorporate real-time scoring and feedback that dynamically updates as users make adjustments to their resumes, helping them optimize documents iteratively.
\item \textbf{Secure Data Handling:} Establish secure data storage protocols and encryption to protect users' personal information and uploaded documents, ensuring compliance with privacy standards and data protection best practices.

\item \textbf{Testing and Iteration:} Conduct testing to ensure that the system’s scoring model and keyword suggestions accurately reflect ATS requirements across various industries. Continuously improve the platform based on user feedback and evolving ATS algorithms.

\end{itemize}


% Applications
\chapter{Applications}


\begin{enumerate}
    \item \textbf{ Resume Optimization:}CV Analyser provides users with tools to optimize their resumes for ATS compatibility. Users can upload their resumes to receive detailed feedback on formatting, keyword relevance, and structure, ensuring that their applications stand out in the hiring process.

    \item \textbf{ Industry-Specific Keyword Recommendations:} The platform offers personalized keyword suggestions based on job descriptions provided by users. This application allows candidates to tailor their resumes effectively, improving their chances of matching the specific criteria outlined by potential employers.

    \item \textbf{Real-Time Resume Scoring:}Users receive a dynamic scoring metric that evaluates their resumes against ATS standards in real time. This feature empowers candidates to make adjustments and see the impact of their changes instantly, facilitating an iterative approach to resume refinement.
    \item \textbf{Personalized Dashboard for Resume Management:} the platform includes a user-friendly dashboard where individuals can manage multiple versions of their resumes, track analysis history, and monitor feedback. This application helps users stay organized and efficiently prepare different resumes for various job applications.

    \item \textbf{ Integration with Job Application Platforms:}The CV Analyser can facilitate integration with popular job application platforms, allowing users to export their optimized resumes directly to job boards. This application streamlines the application process, making it easier for candidates to submit tailored resumes.\\
  
   
\end{enumerate}
In summary the applications of CV Analyser extend beyond simple resume evaluation; they encompass a comprehensive suite of tools and resources that empower job seekers to optimize their applications effectively. By connecting candidates with tailored feedback, industry insights, and supportive community resources, CV Analyser aims to create a more efficient and accessible marketplace for job seekers in a competitive employment landscape.
% References
\chapter{References}
\begin{enumerate}
    \item Smith, J., \textit{Resume Writing for Dummies}, Wiley, 2018. - A comprehensive guide to crafting effective resumes and cover letters.
    \item Brown, C., \textit{The Science of Resume Scoring}, Journal of Career Development, Vol. 45, No. 3, 2021, pp. 25-38. - An analysis of ATS scoring mechanisms and their implications for job seekers.
    \item Lee, K., \& Patel, R., \textit{User Experience Design for Digital Products}, O'Reilly Media, 2021. - Best practices in UX design to enhance user engagement.
    \item Thompson, P., \textit{Leveraging Technology for Effective Job Searches}, Tech Press, 2022. - Discusses the use of technology and platforms for optimizing job applications.
    \item Lutz, M., \textit{Learning Python}, O'Reilly Media, 2013. - A comprehensive resource for beginners and experienced programmers to learn Python.
\end{enumerate}


% Signature
\vfill
\noindent Mrs. Deeksha M \hfill Dr. Manjunath Kotari \\
(Mini-Project Coordinator) \hfill (HOD-CSE)

\end{document}